\section{The PSM as a fixed-excitation slice}
\label{sec:psm-slice}

Fix \(\gls{sym:iexc}=\bar i_{\mathrm{exc}}\) and define
\begin{equation}
 \gls{sym:psipm}=\gls{sym:lexc}\bar i_{\mathrm{exc}}.
 \label{eq:pm-slice}
\end{equation}
Then \cref{eq:model-torque,eq:model-ud,eq:model-uq} become
\begin{align}
 \gls{sym:torque}&=\gls{sym:kt}\gls{sym:iq}
 [\gls{sym:psipm}+\gls{sym:dl}\gls{sym:id}],
 \label{eq:psm-torque}\\
 \gls{sym:ud}&=\gls{sym:rs}\gls{sym:id}
 -\gls{sym:omega}\gls{sym:lq}\gls{sym:iq},\\
 \gls{sym:uq}&=\gls{sym:rs}\gls{sym:iq}
 +\gls{sym:omega}(\gls{sym:ld}\gls{sym:id}+\gls{sym:psipm}).
 \label{eq:psm-voltage}
\end{align}
These are the linear PSM torque and voltage equations. The identification is
geometric, not ontological: a true PSM has no rotor excitation winding and
therefore no term \(\gls{sym:rexc}\bar i_{\mathrm{exc}}^2\). In an EESM slice
that term is a constant loss offset with respect to stator-current allocation.
It must be removed, not silently retained, when interpreting the slice as an
actual PSM.

Set
\begin{equation}
 \gls{sym:ip}=\begin{bmatrix}\gls{sym:id}&\gls{sym:iq}\end{bmatrix}^{\mathsf T},
 \quad
 \gls{sym:mp}=\begin{bmatrix}
 \gls{sym:rs}&-\gls{sym:omega}\gls{sym:lq}\\
 \gls{sym:omega}\gls{sym:ld}&\gls{sym:rs}
 \end{bmatrix},\quad
 \vect r=\begin{bmatrix}0\\\gls{sym:omega}\gls{sym:psipm}\end{bmatrix}.
 \label{eq:psm-affine-map}
\end{equation}
The PSM voltage map is \(\gls{sym:u}=\gls{sym:mp}\gls{sym:ip}+\vect r\).
When \(R_s^2+\omega_e^2L_dL_q>0\), the matrix is invertible, so the inverse
image of the inverter voltage disk is an ellipse. In the EESM picture, every
constant-excitation plane cuts the voltage cylinder in one such ellipse. Its
shape is independent of the chosen excitation; its center moves along the
cylinder null line.

\begin{figure}[tb]
\centering
\begin{tikzpicture}[x={(1cm,0cm)},y={(.32cm,.25cm)},z={(0cm,1cm)},scale=.95]
 \draw[->] (-3,0,0)--(3.1,0,0) node[right]{$i_d$};
 \draw[->] (0,-2,0)--(0,2.1,0) node[right]{$i_q$};
 \draw[->] (0,0,-1.8)--(0,0,2) node[above]{$i_{\mathrm{exc}}$};
 \foreach \z/\x/\col in {-1.2/.75/BrickRed,0/0/ForestGreen,1.2/-.75/MidnightBlue}{
  \begin{scope}[canvas is xy plane at z=\z]
   \draw[very thick,\col] (\x,0) ellipse (1.25 and .66);
  \end{scope}}
 \draw[neutral axis] (.75,0,-1.2)--(-.75,0,1.2);
 \node[annotation] at (2.3,0,1.2){fixed field};
 \node[annotation] at (2.3,0,0){PSM-like slice};
\end{tikzpicture}
\caption{Constant-excitation planes cut the homogeneous EESM voltage cylinder
into translated PSM-like ellipses. Their centers lie on the cylinder's voltage
null line.}
\label{fig:psm-slices}
\end{figure}


The fixed slice also explains the apparent loss of homogeneity. In three
coordinates, scaling all currents scales the field flux and leaves every model
equation homogeneous. In the PSM plane, the field coordinate is no longer
allowed to scale. The term \(\gls{sym:psipm}\) is therefore a constant affine
offset, torque contours are shifted hyperbolas, and voltage ellipses are
translated away from the origin.

\paragraph{What the slice teaches.}
The PSM is the most legible two-dimensional cross-section of the EESM, but a
cross-section must not be confused with a projection. If excitation is free,
projecting the EESM voltage cylinder onto the \(dq\) plane produces an
unbounded strip because field current can cancel \(u_q\). Fixing excitation
produces a bounded ellipse. The distinction becomes decisive whenever a
numerical plot implicitly fixes or bounds excitation.
